\section*{Introduction}
\addcontentsline{toc}{section}{Introduction}
\label{sec:introduction}

Fourteen-year-old Sewell Setzer III did not encounter a foundation model in a laboratory. He encountered a deployed relationship. In \emph{Garcia v. Character Technologies, Inc.}, his mother's complaint alleged that Character.AI combined persistent personas, memory, access rules, warnings, engagement design, and safety controls in a consumer service available to minors. It alleged that Character Technologies' founders and Google made different contributions of personnel, technology, knowledge, and resources. After months of interaction with a character modeled on a fictional television figure, Sewell died by suicide. His mother brought wrongful-death, negligence, product-defect, failure-to-warn, deceptive-practice, and related claims. At the pleading stage, the district court held that the application was plausibly alleged to be a product and declined, on that record, to classify the generated output as protected speech. The action later settled without a finding of defect, causation, or liability.\lrfootnote{See \casecite[1165--69, 1174--84]{garcia2025}; \casecite[1]{garciaDismissal2026}; \casecite[1]{garciaLien2026}. The complaint's disputed account remains an allegation. This Article uses the published Rule 12 order for its classification and pleading analysis, not as a merits judgment.}

The case exposes a problem that survives every possible merits result. The user sees the interface and the output. The facts ordinary law needs lie behind both: which model and application versions ran; who selected the relational objective, system instructions, memory, age gate, crisis response, and engagement metric; what evaluations, complaints, and incident signals arrived; who could update, restrict, roll back, or stop the deployment; and which enterprise retained each record. In a direct-to-consumer service, the operating company may be obvious while its deployment record remains inaccessible. In a divided system, the visible hospital, employer, agency, or application can depend on a model vendor, scoring service, infrastructure provider, and systems integrator that retain different controls. The claimant needs the configuration and control record to test a recognized claim, yet often must identify the record holder before obtaining that record. The merits standard exists; access to the facts that feed it does not.

This Article calls that sequence the \term{pre-merits attribution gap}. It comprises two related failures. The first is an \term{actor-identification failure}: an affected person cannot tell which enterprise selected or retained a control relevant to the asserted risk. The second is a \term{deployment-record failure}: even when the operator's name is visible, the versions, configuration, evaluations, warnings, incidents, and intervention authority needed to test its conduct remain internal or were never recorded. The first failure is most acute in divided deployments. The second can determine a case in any deployment.

Much of the leading AI-liability debate asks a different question: which substantive standard should govern after the relevant actor and conduct are known? Existing negligence, products, antidiscrimination, consumer-protection, administrative, contract, and no-liability rules may answer many such disputes.\lrfootnote{See \artcite{abrahamSharkey2027}; \artcite{geistfeldInsurability2026}; \artcite{hendersonRobotLiability2026}.} This Article accepts that ordinary-law premise and asks what happens one step earlier. A reasonableness standard cannot evaluate conduct until the relevant conduct and actor are visible. Product classification cannot identify which company selected the model harness or retained the update right. A rule denying relief for a category of loss still requires enough information to determine whether the case falls within that category. The front-end information problem is analytically distinct from the standard that ultimately governs the merits.

The question presented is therefore narrow: when a civil claimant can allege a legally cognizable injury or an existing entitlement to equitable relief connected to a designated high-impact AI deployment, but cannot identify the enterprise that controlled the relevant risk or obtain the deployment record needed to test an ordinary cause of action, what threshold should trigger identification, preservation, and focused production---and which enterprise must answer first?

The Article's contribution is institutional sequencing. Existing AI proposals often begin with a substantive duty, disclosure mandate, regulator, or allocation of liability. Existing procedure offers discovery after a viable action reaches an identifiable target. The front door joins four functions that those approaches leave separate: a sectoral designation that identifies the legal interest; an ex ante minimum record that keeps the deployment observable; a visible-operator gateway that makes service possible; and a material-control test that allocates access without allocating liability. The resulting rule is narrow enough to administer and complete enough to open the merits process.

This Article proposes a \term{deployment front-door rule}. A legislature first designates a deployment category and specifies the protected interests, visible operator, minimum record, retention period, trigger, and remedy. Once a claimant alleges a legally cognizable injury or existing equitable entitlement and specific facts supporting a plausible, risk-specific nexus to that deployment, the enterprise that made the deployment available to the affected person or incorporated it into the decision process becomes the disclosure gateway. It must identify the deployed system and relevant entities, preserve the corresponding record, notify enterprises that may hold a material control lane, and provide staged, proportionate, nonprivileged production. Each identified enterprise joins that obligation only if, at the relevant time, it possessed practical authority materially connected to the asserted risk.

The rule changes access to facts, not the elements that decide the case. Its merits firewall creates no presumption of duty, breach, defect, causation, damages, or joint liability. Jurisdiction, immunity, privilege, proportionality, trade-secret protection, and substantive defenses remain operative. An enterprise that completes the first-answer obligation and establishes that it lacked material control over the risk terminates its special front-door role; absent an independently sufficient claim, it returns to ordinary nonparty status. Ultimate allocation remains governed by the cause of action, contribution and indemnity law, contract, and any independently available individual claim.

The front door has an ex ante component because litigation cannot recover a record that a deployment never created. For each designated category, the statute should impose a \term{deployment-record duty}: identify the system and version; assign risk and approval authority; document material evaluations, unresolved risk decisions, changes, complaints, incidents, and interventions; identify who can restrict, roll back, or stop operation; and retain the record for a specified period. What this Article previously called a responsibility charter is best understood as the internal instrument that satisfies that record duty. It is infrastructure for the front door, not a second governance product.

The scope of the required record follows practical control. This Article uses a compressed \term{deployment-control cascade} to map objectives, resources, configuration, feedback, and stopping power across the model provider, application deployer, and operating institution. The cascade is a record-specification method. It identifies evidence relevant to authorization, notice, foreseeability, feasible precaution, causal sequence, preservation, and defense; governing law supplies every legal standard. Token-level unpredictability may complicate causal proof, but it does not make organizational choices about purpose, users, data, instructions, tools, monitoring, updates, and shutdown disappear.

Nor does the proposal depend on resolving artificial personhood. In the deployments covered here, existing people and organizations hold the assets, records, authority, service capacity, and exposure to remedies. A legislature may someday construct a bounded juridical platform for an artificial system. Until it does, operational autonomy remains a fact about how the system acts, while positive law determines which recognized person holds the claim, owes the duty, or answers the judgment. The front door points to the entities that can actually identify the deployment and produce its record.

Three objections define the rule's specifications. First, ordinary pleading and discovery often work once a claimant has a proper defendant and an identifiable record holder. The gap lies in establishing those predicates and in creating the record before litigation. Second, identification and production impose real costs even when the enterprise ultimately prevails. Coverage must therefore be sectorally designated; the nexus must be risk specific; production must be staged; and privacy, security, privilege, and trade secrets must be protected. Third, some losses connected to AI have no legal remedy. The rule requires an existing cause of action or equitable entitlement, permits pure legal defenses to be resolved first, and supplies no substantive right to inspect a deployment simply because a claimant invokes AI.

The Article develops one contribution through five Parts. Part I identifies the pre-merits attribution gap, explains why ordinary procedure presupposes facts it cannot always reveal, and compresses the deployment-control cascade into a tool for specifying the missing record. Part II states the front-door rule, defines sectoral coverage and material control, separates record creation from preservation and production, supplies model statutory text, and coordinates state legislation with federal procedure. Part III applies the rule to two complementary settings: \emph{Garcia}, where the visible operator makes production the central problem, and \emph{Mobley}, where hiring decisions, platform features, data hosting, legal control, and validation records divide across employers and a vendor. A completed public-employment record provides a stable analogue. Part IV develops the ex ante record infrastructure and shows how existing safety and information-demand regimes make it administrable. Part V defines the merits firewall, information-cost controls, fragmented-stack boundaries, and classification neutrality. The Conclusion returns to a single proposition: substantive law should decide AI cases, not an information asymmetry that prevents substantive law from starting.
